% Generated by scripts/assemble.py; corrections belong in source/corrections.json.
\examhead{54}{华东理工大学 2026 年数学分析试卷}{55}
% p056-b01; source PDF page 56
\par\Needspace{5\baselineskip}
\blockanchor{p056-b01}
\majorquestion{一}解答如下问题:\par
% p056-b02; source PDF page 56
\blockanchor{p056-b02}
\question{1}写出基本数列(Cauchy 数列)的定义.\par
% p056-b03; source PDF page 56
\blockanchor{p056-b03}
\question{2}设 $x_{1} > 0$, 对 $n \geq 1$, 有 $x_{n+1} = \arctan x_{n}$, 证明: $\{x_{n}\}$ 是基本数列.\par
% p056-b04; source PDF page 56
\par\Needspace{5\baselineskip}
\par\Needspace{6\baselineskip}
\blockanchor{p056-b04}
\majorquestion{二}设函数 $f$ 在闭区间 $[a, b]$ 上有二阶导数，且 $f'(a) = f'(b) = 0$，证明：存在 $\xi \in (a, b)$，使得\par
\nopagebreak[4]
% p056-b05; source PDF page 56
\blockanchor{p056-b05}
\displayformula{| f ^ {\prime \prime} (\xi) | \geq 4 \frac {| f (b) - f (a) |}{(b - a) ^ {2}}.}
% p056-b06; source PDF page 56
\par\Needspace{5\baselineskip}
\par\Needspace{6\baselineskip}
\blockanchor{p056-b06}
\majorquestion{三}在 $\mathbb{R}^2$ 上的函数\par
\nopagebreak[4]
% p056-b07; source PDF page 56
\blockanchor{p056-b07}
\displayformula{f (x, y) = \left\{ \begin{array}{l l} \frac {x ^ {2} - y ^ {2}}{x ^ {2} + y ^ {2}}, & x ^ {2} + y ^ {2} \neq 0, \\ 0, & x ^ {2} + y ^ {2} = 0. \end{array} \right.}
% p056-b08; source PDF page 56
\blockanchor{p056-b08}
讨论 f 在  $(0,0)$  点的连续性.\par
% p056-b09; source PDF page 56
\par\Needspace{5\baselineskip}
\blockanchor{p056-b09}
\majorquestion{四}解答如下问题:\par
% p056-b10; source PDF page 56
\blockanchor{p056-b10}
\question{1}$f$ 在 $[0,1]$ 上可积, 说明: $\int_0^1 f(x)\mathrm{d}x = \lim_{n\to \infty}\sum_{k = 1}^{n}\frac{1}{n} f\left(\frac{k - 1}{n}\right)$.\par
% p056-b11; source PDF page 56
\blockanchor{p056-b11}
\question{2}对 $n \geq 1$, 定义 $A_{n} = \frac{\sum_{k=1}^{n} k \sin\left(\frac{k-1}{n}\pi\right)}{n^{2}}$, 求 $\lim_{n \to \infty} A_{n}$.\par
% p056-b12; source PDF page 56
\par\Needspace{5\baselineskip}
\blockanchor{p056-b12}
\majorquestion{五}记 $I_{p} = \int_{2}^{+\infty} \frac{(\ln x)^{p} \sin x}{x} \, \mathrm{d}x, p \in \mathbb{R}$.\par
% p056-b13; source PDF page 56
\blockanchor{p056-b13}
\question{1}当 $p = 0$ 时, 求 $I_{p}$ 的敛散性.\par
% p056-b14; source PDF page 56
\blockanchor{p056-b14}
\question{2}当 $p = 1$ 时, 求 $I_{p}$ 的敛散性.\par
% p056-b15; source PDF page 56
\par\Needspace{5\baselineskip}
\par\Needspace{6\baselineskip}
\blockanchor{p056-b15}
\majorquestion{六}计算曲线积分 $I = \oint_{\Gamma} |\sqrt{3}y - x| \, \mathrm{d}x - 5z \, \mathrm{d}z$，其中\par
\nopagebreak[4]
% p056-b16; source PDF page 56
\blockanchor{p056-b16}
\displayformula{\Gamma : \left\{ \begin{array}{l} x ^ {2} + y ^ {2} + z ^ {2} = 8, \\ x ^ {2} + y ^ {2} = 2 z. \end{array} \right.}
% p056-b17; source PDF page 56
\blockanchor{p056-b17}
从 z 轴正向往坐标原点看去取逆时针方向.\par
% p056-b18; source PDF page 56
\par\Needspace{5\baselineskip}
\par\Needspace{6\baselineskip}
\blockanchor{p056-b18}
\majorquestion{七}已知 $F$ 有二阶连续偏导数, $z = z(x, y)$ 是由 $F\left(z + \frac{1}{x}, z - \frac{1}{y}\right) = 0$ 确定的隐函数. 且 $z$ 有连续二阶偏导数. 证明:  

\question{1}$x^{2} \frac{\partial z}{\partial x} - y^{2} \frac{\partial z}{\partial y} = 1$.\par
\nopagebreak[4]
% p056-b19; source PDF page 56
\blockanchor{p056-b19}
\displayformula{x ^ {3} \frac {\partial^ {2} z}{\partial x ^ {2}} + x y (x - y) \frac {\partial^ {2} z}{\partial x \partial y} - y ^ {3} \frac {\partial^ {2} z}{\partial y ^ {2}} + 2 = 0.}
% p056-b20; source PDF page 56
\par\Needspace{5\baselineskip}
\blockanchor{p056-b20}
\majorquestion{八}设 $\Sigma$ 是单位球面 $\{(x, y, z) \in \mathbb{R}^3 : x^2 + y^2 + z^2 = 1\}$, 记 $I_p = \iint_{\Sigma} (x + y + z)^p \mathrm{d}S$, $p$ 是正整数, 求 $I_1, I_2$.
\majorquestion{九}解答如下问题:\par
% p056-b21; source PDF page 56
\blockanchor{p056-b21}
\question{1}给出函数列一致收敛的定义.\par
% p056-b22; source PDF page 56
\blockanchor{p056-b22}
\question{2}对 $n \geq 1$, 设 $S_{n}(x) = \frac{x}{1 + n^{2}x^{2}}$, 证明: 函数列 $\{S_{n}(x)\}$ 在 $(0, +\infty)$ 上一致收敛.\par
% p056-b23; source PDF page 56
\blockanchor{p056-b23}
\question{3}当 $n \geq 1$ 时, $T_{n}(x) = \frac{\sqrt{x}}{1 + n^{2}x^{2}}$, 问: 函数列 $\{T_{n}(x)\}$ 在 $(0, +\infty)$ 上是否一致收敛? 请证明你的结论.\par
